\begin{figure}[!ht]
\centering
\resizebox{\textwidth}{!}{%
\begin{tikzpicture}[
    font=\small,
    backend/.style={draw=subsectionblue, fill=subsectionblue!10, rounded corners=4pt, align=center, minimum width=4cm, minimum height=1.6cm, thick},
    llm/.style={draw=ugrred, fill=ugrred!10, rounded corners=4pt, align=center, minimum width=4cm, minimum height=1.6cm, thick},
    front/.style={draw=orange!90, fill=orange!10, rounded corners=4pt, align=center, minimum width=4cm, minimum height=1.6cm, thick},
    arrow/.style={-{Stealth[length=2.5mm]}, thick, draw=gray!80!black},
    lbl/.style={font=\scriptsize, inner sep=4pt, text=gray!80!black}
]

% 1. Nodos mucho más separados (en X y en Y)
\node[backend] (sql)   at (0, 0) {\textbf{Base de Datos (Backend)}\\Resultados SQL completos};
\node[llm]     (llm)   at (7.5, 3) {\textbf{Modelo de Lenguaje}\\Diseño del componente visual};
\node[backend] (merge) at (7.5, -3) {\textbf{Motor Visual (Backend)}\\Inyección de datos};
\node[front]   (ui)    at (15.5, -3) {\textbf{Interfaz (Frontend)}\\Renderizado visual}; % Nodo desplazado más a la derecha

% 2. Flecha hacia el LLM (Arriba, giro suave)
\draw[arrow, rounded corners=8pt] (sql.north) |- node[lbl, above=2pt, pos=0.75] {1. Metadatos y muestra (3 filas)} (llm.west);

% 3. Flecha del LLM al Motor Visual (Centro, recta)
\draw[arrow] (llm.south) -- node[lbl, right=2pt] {2. Esqueleto JSON vacío} (merge.north);

% 4. Flecha de la BD al Motor Visual (Abajo, giro suave)
\draw[arrow, rounded corners=8pt] (sql.south) |- node[lbl, below=2pt, pos=0.75] {3. Inserción del 100\% de las filas} (merge.west);

% 5. Flecha al Frontend (Derecha, recta con espacio amplio)
\draw[arrow] (merge.east) -- node[lbl, above=2pt] {4. Envío por WebSocket} (ui.west);

\end{tikzpicture}%
}
\caption{Arquitectura del Motor Visual Generativo (\textit{Generative UI}).} %El modelo de lenguaje actúa únicamente como diseñador de la estructura visual, mientras que el \emph{backend} asegura la integridad de la información inyectando los datos reales. Esto ahorra \textit{tokens} y elimina el riesgo de alucinaciones numéricas.}
\label{fig:generative_ui}
\end{figure}